\documentclass[12pt]{exam}
%\documentclass[11pt,a4paper]{exam}
\usepackage{amsmath,amsthm,amsfonts,amssymb,dsfont}
\usepackage{float}
\usepackage{ifthen}
\usepackage{array}
\usepackage{geometry}
\geometry{
legalpaper, total={177.8mm, 290mm},left=10mm, right=10mm,
top=7mm, bottom=10mm,
}
\usepackage{enumerate}% http://ctan.org/pkg/enumerate
\usepackage{multicol}
\usepackage{hhline}
\usepackage[table]{xcolor}


% Accumulate the answers. Unmodified from Phil Hirschorn's answer
% https://tex.stackexchange.com/questions/15350/showing-solutions-of-the-questions-separately/15353
\newbox\allanswers
\setbox\allanswers=\vbox{}

\newenvironment{answer}
{%
    \global\setbox\allanswers=\vbox\bgroup
    \unvbox\allanswers
}%
{%
    \bigbreak
    \egroup
}

\newcommand{\showallanswers}{\par\unvbox\allanswers}
% End Phil's answer


% Is there a better way?
\newcommand*{\getanswer}[5]{%
    \ifthenelse{\equal{#5}{a}}
    {\begin{answer}\thequestion. (a)~#1\end{answer}}
    {\ifthenelse{\equal{#5}{b}}
        {\begin{answer}\thequestion. (b)~#2\end{answer}}
        {\ifthenelse{\equal{#5}{c}}
            {\begin{answer}\thequestion. (c)~#3\end{answer}}
            {\ifthenelse{\equal{#5}{d}}
                {\begin{answer}\thequestion. (d)~#4\end{answer}}
                {\begin{answer}\textbf{\thequestion. (#5)~Invalid answer choice.}\end{answer}}}}}
}

\setlength\parindent{0pt}
%usage \choice{ }{ }{ }{ }
%(A)(B)(C)(D)
\newcommand{\fourch}[5]{
    \par
    \begin{tabular}{*{4}{@{}p{0.23\textwidth}}}
        (a)~#1 & (b)~#2 & (c)~#3 & (d)~#4
    \end{tabular}
    \getanswer{#1}{#2}{#3}{#4}{#5}
}

%(A)(B)
%(C)(D)
\newcommand{\twoch}[5]{
    \par
    \begin{tabular}{*{2}{@{}p{0.46\textwidth}}}
        (a)~#1 & (b)~#2
    \end{tabular}
    \par
    \begin{tabular}{*{2}{@{}p{0.46\textwidth}}}
        (c)~#3 & (d)~#4
    \end{tabular}
    \getanswer{#1}{#2}{#3}{#4}{#5}
}

%(A)
%(B)
%(C)
%(D)
\newcommand{\onech}[5]{
    \par
    (a)~#1 \par (b)~#2 \par (c)~#3 \par (d)~#4
    \getanswer{#1}{#2}{#3}{#4}{#5}
}

\newlength\widthcha
\newlength\widthchb
\newlength\widthchc
\newlength\widthchd
\newlength\widthch
\newlength\tabmaxwidth

\setlength\tabmaxwidth{0.96\textwidth}
\newlength\fourthtabwidth
\setlength\fourthtabwidth{0.25\textwidth}
\newlength\halftabwidth
\setlength\halftabwidth{0.5\textwidth}

\newcommand{\choice}[5]{%
\settowidth\widthcha{AM.#1}\setlength{\widthch}{\widthcha}%
\settowidth\widthchb{BM.#2}%
\ifdim\widthch<\widthchb\relax\setlength{\widthch}{\widthchb}\fi%
    \settowidth\widthchb{CM.#3}%
\ifdim\widthch<\widthchb\relax\setlength{\widthch}{\widthchb}\fi%
    \settowidth\widthchb{DM.#4}%
\ifdim\widthch<\widthchb\relax\setlength{\widthch}{\widthchb}\fi%

% These if statements were bypassing the \onech option.
% \ifdim\widthch<\fourthtabwidth
%     \fourch{#1}{#2}{#3}{#4}{#5}
% \else\ifdim\widthch<\halftabwidth
% \ifdim\widthch>\fourthtabwidth
%     \twoch{#1}{#2}{#3}{#4}{#5}
% \else
%      \onech{#1}{#2}{#3}{#4}{#5}
%  \fi\fi\fi}

% Allows for the \onech option.
\ifdim\widthch>\halftabwidth
    \onech{#1}{#2}{#3}{#4}{#5}
\else\ifdim\widthch<\halftabwidth
\ifdim\widthch>\fourthtabwidth
    \twoch{#1}{#2}{#3}{#4}{#5}
\else
    \fourch{#1}{#2}{#3}{#4}{#5}
\fi\fi\fi}


\begin{document}



\begin{table}[]
\begin{tabular}{llllllllcllll}
\textit{Deneb} &  &  &  &  &  &  &  & \multicolumn{1}{l}{\textbf{MYMENSINGH GIRLS' CADET COLLEGE}} &                                             &                        &                                 &                        \\
                &  &  &  &  &  &  &  & \multicolumn{1}{c}{PROGRESS TEST EXAMINATION - 2025}       &                                             &                        & \multicolumn{1}{c}{}            &                        \\
                &  &  &  &  &  &  &  & CLASS: HSC                                                   &                                             &                        & \multicolumn{1}{c}{}            &                        \\
                &  &  &  &  &  &  &  & MULTIPLE CHOICE QUESTIONS                                    &                                             &                        & \multicolumn{1}{c}{}            &                        \\
                &  &  &  &  &  &  &  & STATISTICS                                                   &                                             &                        & \multicolumn{1}{r}{}            &                        \\ \cline{11-13} 
                &  &  &  &  &  &  &  & FIRST PAPER                                                 & \multicolumn{1}{r|}{\textbf{Subject Code:}} & \multicolumn{1}{l|}{1} & \multicolumn{1}{l|}{\textbf{2}} & \multicolumn{1}{l|}{9} \\ \cline{11-13} 
                &  &  &  &  &  &  &  & [According to the Syllabus of 2026]                          & \multicolumn{1}{r}{}                        &                        &                                 &                        \\ \cline{12-12}
                &  &  &  &  &  &  &  & TIME – 20 minutes                                            & \multicolumn{1}{r}{\textbf{Set:}}           & \multicolumn{1}{l|}{}  & \multicolumn{1}{l|}{\textbf{Ka}} &                        \\ \cline{12-12}
                &  &  &  &  &  &  &  & FULL MARKS – 20                                              &                                             &                        &                                 &                       
\end{tabular}
\end{table}

%  \normalfont\normalsize
 % 11.45a.m.~--~1.45p.m.
\hrule

\begin{center}
[N.B. – Answer all the questions. Each question carries ONE mark. Block fully, with a black ball- point pen, the circle of the letter that stands for the correct/best answer in the “Answer sheet” for the Multiple Choice Questions Examination.]\\

  
  \textbf{Candidates are asked not to leave any mark or spot on the question paper.}
\end{center}
\begin{questions}

\question \textbf{A survey categorizing people by their favorite color is an example of which measurement scale?}  
\choice{Nominal}{Ordinal}{Interval}{Ratio}{a} 

\question \textbf{Which is a discrete variable?}
\choice{Number of languages spoken by a person}
{Time taken to complete a race}{Length of a road}{Volume of water in a tank}{a}

\question \textbf{If $x_1=3$, $x_2=2$, $x_3=-6$, and $x_4=4$, what is
$\displaystyle \sum_{i=1}^4 x_i^2$?}  
\choice{45}{65}{85}{89}{b}

% Situation Set 1 Starts
\textbf{Answer the next TWO questions based on the following information}

The table below shows monthly sales (in thousands of dollars) for two products in four regions.

\begin{table}[h]
\centering
\begin{tabular}{c|cccc}
Product & North & South & East & West \\ \hline
A & 12 & 15 & 10 & 8 \\
B & 18 & 22 & 14 & 16
\end{tabular}
\end{table}

\question \textbf{Calculate total sales of Product A across all regions.}
\choice{40}{42}{45}{48}{c}

\question \textbf{Calculate total sales of all products across all regions.}
\choice{95}{100}{105}{115}{d}
% Situation Set 1 Ends

\end{questions}

\vspace{1cm}
\vfill
\begin{center}
%“You can have data without information, but you cannot have information without data.” \\ %- Daniel Keys Moran
--AAM-- 
\end{center}

\pagebreak
%\newpage  %Uncomment to put on new age
\bigskip

%\begin{multicols}{3}
%[
%Answer Key
%]
%\showallanswers % Phil Hirschorn
%\end{multicols}


\end{document}